\documentclass[10pt,journal,compsoc]{IEEEtran}

\usepackage{amsmath,amssymb}
\usepackage{booktabs}
\usepackage{tabularx}
\usepackage{array}
\usepackage{url}
\usepackage{microtype}

\newcolumntype{Y}{>{\raggedright\arraybackslash}X}

\title{Supplementary Material for ``Information-Set Conditional Integrity Auditing for Hybrid Quantum-Classical Kernel Workflows''}
\author{Roberto~Fern\'andez-Barrios, Iker~Pastor-L\'opez, Amaia~Pikatza-Huerga,
and~Pablo~Garc\'ia~Bringas\\ Faculty of Engineering, University of Deusto,
Bilbao, Spain}

\begin{document}
\maketitle

\section{Purpose and Claim Boundary}

This supplement records the complete experimental design, coverage contract,
validation checks, and reproduction mapping behind the main article. It adds no
QPU, calibrated-noise, scheduling, provider-security, multi-tenancy, or
operational Fleet Management evidence. The exact claims remain bounded to
data-to-evaluation interventions, ideal-statevector quantum kernels, controlled
binomial shot emulation, and a local executable contract.

\section{Frozen Experimental Design}

\subsection{Expansion Environments}

The expansion comprises eight fixed environments (Table~\ref{tab:env}). ``OOD''
means a prespecified source-to-target construction, not a sampled deployment
population. Every environment uses five split seeds and two nested model seeds.
All evaluate projected dimensions 8, 10, and 12. The CICIDS scale gate uses
256/256 train/evaluation samples; all others use 128/128. The two additional ID
datasets retain all empirically distinct quantum profiles. Scale and temporal
gates compare SVC-RBF with QSVC-ZZ.

\begin{table*}[!t]
\caption{Prespecified expansion environments.}
\label{tab:env}
\centering
\footnotesize
\begin{tabularx}{\textwidth}{@{}p{0.05\textwidth}p{0.24\textwidth}p{0.20\textwidth}p{0.20\textwidth}YY@{}}
\toprule
ID & Environment & Training source & Evaluation source & Samples & Profiles \\
\midrule
E1 & CICIDS scale/ID & Frozen CICIDS subset & Same staged source & 256/256 & SVC, ZZ \\
E2 & UNSW-NB15 ID & Balanced staged UNSW & Same staged source & 128/128 & SVC, ZZ, PauliXYZ, Z/PauliXZ-equivalent \\
E3 & ToN-IoT ID & Balanced staged ToN-IoT & Same staged source & 128/128 & SVC, ZZ, PauliXYZ, Z/PauliXZ-equivalent \\
E4 & CICIDS Tue$\rightarrow$Wed & Tuesday traffic & Wednesday traffic & 128/128 & SVC, ZZ \\
E5 & CICIDS Tue$\rightarrow$Fri-PortScan & Tuesday traffic & Friday PortScan traffic & 128/128 & SVC, ZZ \\
E6 & CICIDS Wed$\rightarrow$Thu-Web & Wednesday traffic & Thursday Web traffic & 128/128 & SVC, ZZ \\
E7 & CICIDS Wed$\rightarrow$Fri-AM & Wednesday traffic & Friday morning traffic & 128/128 & SVC, ZZ \\
E8 & UNSW temporal OOD & Prespecified UNSW source period & Prespecified UNSW target period & 128/128 & SVC, ZZ \\
\bottomrule
\end{tabularx}
\end{table*}

Each staged table contains 3,000 balanced binary rows before experimental
subsampling. Numeric non-finite values are replaced during staging; categorical
fields are excluded. These choices produce a controlled integrity audit but do
not reproduce operational prevalence. Several temporal targets yield clean
performance near chance, limiting the impact that later interventions can
express; the exact information-set invariances do not depend on clean headroom.

\subsection{Perturbation Suite}

The paper-core suite crosses six mechanisms with nominal strengths 0.02, 0.05,
and 0.10 (Table~\ref{tab:attacks}). Artifact identifiers retain the historical
family string \texttt{target\_shift}; the scientific text uses
\emph{evaluation-label intervention} because these rows do not instantiate the
statistical label-shift problem.

\begin{table}[!t]
\caption{Frozen 18-condition intervention suite.}
\label{tab:attacks}
\centering
\footnotesize
\begin{tabularx}{\columnwidth}{@{}p{0.08\columnwidth}Yp{0.26\columnwidth}@{}}
\toprule
IDs & Mechanism & Strengths \\
\midrule
1--3 & Feature sign flip & 0.02, 0.05, 0.10 \\
4--6 & Feature-wise mean shift & 0.02, 0.05, 0.10 \\
7--9 & Scaling drift & 0.02, 0.05, 0.10 \\
10--12 & Feature dropout & 0.02, 0.05, 0.10 \\
13--15 & Prior-preserving evaluation-label flip & 0.02, 0.05, 0.10 \\
16--18 & Random evaluation-label flip & 0.02, 0.05, 0.10 \\
\bottomrule
\end{tabularx}
\end{table}

The feature mechanisms may alter representation, scores, predictions, and the
reported conclusion. Both label mechanisms keep the feature representation,
scores, and predictions fixed. Prior-preserving flips also retain the empirical
class histogram. Equal mechanism weighting is a deliberately transparent stress
profile, not an estimate of real incident frequency.

\section{Information-Set Coverage Contract}

Table~\ref{tab:regimes} records what each evidence regime can establish in the
implemented study. Structural blindness means that the sensor's complete input
is unchanged by construction; it is stronger and narrower than observing zero
empirical response in a finite sample.

\begin{table*}[!t]
\caption{Evidence regimes, implemented sensors, and demonstrated blind regions.}
\label{tab:regimes}
\centering
\footnotesize
\begin{tabularx}{\textwidth}{@{}p{0.10\textwidth}p{0.19\textwidth}p{0.28\textwidth}YY@{}}
\toprule
Regime & Evidence & Implemented sensor family & Exact/empirical blind region & Minimum closing evidence in this study \\
\midrule
$\mathcal I_X$ & Features & Feature JSD, MMD, KS rejection & Every label-only change is structurally blind & Trusted label or item provenance \\
$\mathcal I_{XF}$ & Features, scores, predictions & $\mathcal I_X$ plus score JSD and prediction disagreement & Every label-only change is structurally blind for a fixed predictor & Item-aligned outcome-prediction evidence \\
$\mathcal I_{Y_m}$ & Label class counts/marginal & Prior delta and label JSD & Prior-preserving item-level relabeling is structurally blind & Trusted item-level label provenance \\
$\mathcal I_{XFY}$ & Item-aligned features, predictions, outcomes & Prediction-rate, disagreement, prediction-distribution, confusion-profile deltas & Closes every positive label-path case executed, conditional on trusted outcomes & Independent authentication of the joint evidence \\
$\mathcal I_Q$ & Circuit, parameters, kernel, estimation evidence & Hashes, semantic probes, algebra, repeated estimation, output changes & Sensor-dependent: hashes overreact to equivalence; algebra misses PSD-preserving substitution; outputs miss sub-decision changes & Layered policy-aware composition \\
\bottomrule
\end{tabularx}
\end{table*}

Every machine-readable coverage record contains seven fields:
\begin{enumerate}
  \item workflow boundary;
  \item protected asset;
  \item adversary or failure capability, including fixed elements;
  \item available information regime;
  \item sensor, reference, threshold, and false-alarm budget;
  \item covered perturbations and residual blind region; and
  \item response or countermeasure.
\end{enumerate}
This record prevents a sensor response from being promoted into a global
security claim.

\section{Deduplication and Statistical Units}

Gate 1 emits 7,980 raw model-intervention rows. Different quantum-map jobs
repeat the same classical SVC pipeline, producing 3,420 repeated rows. The
builder compares every repeated key across performance, impact, and raw signal
columns before retaining one representative; 4,560 unique observations remain.

The expansion emits 13,680 raw rows from 360 completed CSV/JSON job pairs. Its
multi-map ID gates contain 2,280 repeated SVC rows, again accepted as duplicates
only after numerical identity checks; 11,400 unique observations remain. The
strict builder fails if a job is missing, an environment contains an unexpected
model/dimension/seed/attack design, or a duplicated row disagrees.

Coverage counts use prespecified cells as their denominators. For the secondary
suite-average model profile, nested model-seed values are averaged inside each
of five split seeds. The split seed is the inferential unit, and the interval is
\begin{equation}
 \bar d \pm t_{0.975,4}\frac{s_d}{\sqrt{5}},
\end{equation}
where $d$ is the within-split ZZ-minus-SVC impact difference. The intervals are
not multiplicity-adjusted and do not support population transport. Results over
20 split-model combinations are sensitivity analyses only.

\section{Exact-Statevector Engine Validation}

The expansion evaluator replaces pairwise compute-uncompute execution with a
mathematically equivalent ideal-statevector calculation. It prepares each
unique state $\lvert\phi(x)\rangle$ once and forms
\begin{equation}
 K(x,y)=\left|\langle\phi(x)\mid\phi(y)\rangle\right|^2
\end{equation}
by matrix multiplication. Symmetric blocks use the same eigenvalue-clipping PSD
repair as the frozen Qiskit reference.

Validation has four layers:
\begin{enumerate}
  \item symmetric and cross-kernel unit matrices agree with the reference at
  relative and absolute tolerance $10^{-10}$;
  \item symmetry, unit diagonal, numerical range, cache bounds, and cache reuse
  are unit-tested;
  \item a CICIDS 128/128 end-to-end replay (split/model seed 42, dimension 8,
  ZZ, paper-core suite) preserves every categorical, discrete performance,
  impact, and integrity-signal endpoint; and
  \item the prespecified 256/256 pilot completes.
\end{enumerate}

The two end-to-end outputs each have 38 rows and 126 columns. Only ROC-AUC and
its derived drops differ, with maximum absolute difference
$1.221896\times10^{-4}$, approximately one half-tie increment for 64 positive
and 64 negative evaluation examples. This supports an explicit $2\times10^{-4}$
ROC-AUC tolerance, not bitwise identity of continuous scores. The engine is
tagged \texttt{exact\_statevector}; it supplies no finite-shot, noise-model, or
hardware evidence.

\section{Quantum Gate: Conditions and Coverage}

The quantum gate crosses five split seeds, dimensions 4/6/8, and 11 conditions,
giving 165 cells. Table~\ref{tab:qcoverage} reports response fractions over the
15 split-dimension cells per condition. A value of one means nonzero raw
response in all 15 cells at the frozen tolerance; it is not calibrated power.

\begin{table*}[!t]
\caption{Quantum-gate sensor response fractions.}
\label{tab:qcoverage}
\centering
\scriptsize
\setlength{\tabcolsep}{3.5pt}
\begin{tabular}{@{}lrrrrrr@{}}
\toprule
Condition & Circuit prov. & Kernel prov. & Semantic kernel & Algebraic kernel & Repeated estimate & Output change \\
\midrule
Clean & 0.00 & 0.00 & 0.00 & 0.00 & 0.00 & 0.00 \\
Benign transpilation & 1.00 & 1.00 & 0.00 & 0.00 & 0.00 & 0.00 \\
Common-unitary rewrite & 1.00 & 1.00 & 0.00 & 0.00 & 0.00 & 0.00 \\
Parameterized RZ 0.02 & 1.00 & 1.00 & 1.00 & 0.00 & 0.00 & 0.00 \\
Parameterized RZ 0.10 & 1.00 & 1.00 & 1.00 & 0.00 & 0.00 & 0.13 \\
Feature-map reps 1$\rightarrow$2 & 1.00 & 1.00 & 1.00 & 0.00 & 0.00 & 0.73 \\
Asymmetric kernel edit & 0.00 & 1.00 & 1.00 & 1.00 & 0.00 & 0.00 \\
Kernel diagonal erosion & 0.00 & 1.00 & 1.00 & 1.00 & 0.00 & 0.00 \\
PSD-preserving kernel mix & 0.00 & 1.00 & 1.00 & 0.00 & 0.00 & 0.07 \\
Shot estimator 256 & 0.00 & 1.00 & 1.00 & 1.00 & 1.00 & 0.47 \\
Shot estimator 1,024 & 0.00 & 1.00 & 1.00 & 1.00 & 1.00 & 0.20 \\
\bottomrule
\end{tabular}
\end{table*}

All nine frozen acceptance checks pass: 165 cells are present; clean evidence is
invariant; benign transpilation and common-unitary rewrite alter provenance but
preserve semantic kernels and outputs; every data-dependent circuit mutation
changes the semantic kernel; symmetry and diagonal sensors catch their
respective malformed edits; the PSD-preserving mixture passes algebra but is
exposed by trusted provenance/semantics; and both shot estimators produce
nonzero repeat discrepancy in every cell.

\section{Executable Contract Details}

The prototype evaluates four ordered contracts:
\begin{enumerate}
  \item \textbf{Input/preprocessing}: input-array and preprocessing-declaration
  hashes; mismatch blocks.
  \item \textbf{Circuit/kernel}: circuit/kernel provenance, semantic probes,
  symmetry, diagonal, eigenvalue, and PSD-repair evidence; an equivalent rewrite
  passes only if explicitly approved.
  \item \textbf{Execution/result}: declared execution metadata, repeated-
  estimation discrepancy, and prediction provenance; approved stochastic
  execution is held for adjudication.
  \item \textbf{Evaluation/report}: item-level labels are bound to predictions
  and balanced accuracy is recomputed; prior-preserving corruption blocks even
  when the marginal is identical.
\end{enumerate}

Contracts form an ordered SHA-256 chain. Later modification is detectable, but
the chain does not authenticate a provider because it is not digitally signed.
The response lattice is deliberately fail-closed: all passes yield
\texttt{allow}; at least one review and no failure yields \texttt{hold}; any
failure yields \texttt{block}. Six prespecified scenarios and all six chains
pass eight strict acceptance checks in the frozen artifact.

\section{Preregistered Reinforcement Gates (Artifact 1.1.0)}

Three sensitivity gates were added after the 1.0.0 evidence was frozen. Their
protocol, endpoints, and one implementation clarification (amendment A1) were
recorded in \path{manuscript/paper15_v11_reinforcement_prereg.md} before
execution; results are reported as executed. The gates reuse the exact-
statevector engine, the frozen split and model seeds, and the frozen
18-condition suite. They add no QPU, calibrated-noise, scheduling, provider,
multi-tenancy, or Fleet Management evidence, and they do not change any 1.0.0
count.

\subsection{Gate N: Null Calibration with Disjoint Clean Pools}

For each of the 240 expansion runs (eight environments, five split seeds, two
model seeds, three dimensions), the clean pool is the set of evaluation-side
rows not selected into the frozen evaluation batch (900\,$-$\,$n$ in ID
protocols, 3{,}000\,$-$\,$n$ in OOD protocols), transformed with the same
training-fitted projection and scaler as the frozen run. A permutation seeded
by the split seed divides the pool into a calibration half and an evaluation
half. Twenty clean batches of $n$ rows are drawn from each half by simple
random sampling without replacement; the suite also includes an identity
control and two near-null controls (Gaussian noise $\sigma=0.001$, scaling
$\alpha=0.001$) applied in place.

Ten distributional sensors are calibrated: feature JSD, MMD, and KS rejection
rate against the clean batch; score JSD; prediction-rate shift and prediction
JSD; label prior shift and label JSD; and the two confusion-profile statistics.
For each (environment, dimension, model) cell the 200 calibration values are
pooled and the threshold is the 191st order statistic
($\lceil 201\times0.95\rceil$, $\alpha=0.05$); a sensor fires when its value
strictly exceeds the threshold. Thresholds are frozen before any attacked row
is examined. The empirical false-alarm rate is the firing fraction on the 200
disjoint evaluation draws of the same cell. A regime fires when any of its
sensors fires; regime-level false-alarm rates are therefore unions and are
reported as such, without correction.

Two auditors are distinguished (amendment A1). The \emph{batch-level auditor}
receives a new batch without item correspondence and can use only the
calibrated distributional sensors. The \emph{item-aligned auditor} compares
the same items against trusted references, so prediction disagreement,
confusion-profile deltas, and item-level label mismatch have an exactly zero
null and any non-zero delta is a detection. Item-aligned sensors are undefined
on clean draws, whose items differ from the reference batch, and are excluded
from batch-level unions. Detection rates apply the frozen thresholds to the
frozen \texttt{paper\_core} observations of the same cell.

Two disclosed limitations are fixed in advance. In E1 ($n=256$) each pool half
holds 322 rows, so successive draws overlap heavily and E1 calibrated rates
are indicative only. Simple random draws let class counts vary, whereas the
frozen batch is stratified, so the label-marginal null reflects natural batch
variability.

\textbf{Results.} All 240 runs completed (14/14 acceptance checks).
Table~\ref{tab:s-fpr} reports the pooled false-alarm rates: every sensor lies
between 0.029 and 0.069 against the nominal 0.05, and the per-cell rates have
median 0.04 with 92.8\% of non-E1 cells inside $[0.01, 0.10]$.
Table~\ref{tab:s-fpr-env} shows the environment dependence; E1 and E3 (ToN-IoT,
25.7\% near-duplicate rows) have the highest regime unions.
Table~\ref{tab:s-detection} gives the calibrated detection rates and the
near-null control responses; Table~\ref{tab:s-labelpath} isolates the label
path: the calibrated $\mathcal I_X$, $\mathcal I_{XF}$ and (prior-preserving)
$\mathcal I_{Y_m}$ regimes fire on none of the 3,600 label observations, the
batch-level $\mathcal I_{XFY}$ union fires on 43 (40 of the 2,184 positive
impacts), and the item-aligned auditor detects all 3,600. Table~\ref{tab:s-mass}
records the cluster structure that makes the KS sensor react to in-place
perturbations of order $10^{-3}$ standard deviations while its between-batch
null stays narrow; this fingerprint is legitimate for unmodified data but would
be evaded by perturbations that preserve the clusters.

\begin{table}[!t]
\caption{Gate N: empirical false-action rate of the null-calibrated sensors and batch-level regimes on 12,000 clean draws from evaluation pools disjoint from the calibration pools (all eight environments pooled; nominal $\alpha=0.05$ per sensor). Resampled draws within a pool may overlap. Intervals are descriptive Clopper--Pearson 95\% intervals: draws within a cell share a pool half and are not independent. Regime rows are unions without correction.}
\label{tab:s-fpr}
\centering
\footnotesize
\setlength{\tabcolsep}{4pt}
\begin{tabular}{@{}lrrl@{}}
\toprule
Sensor / regime & Fires & Rate & 95\% CI \\
\midrule
Feature JSD & 600 & 0.050 & [0.046, 0.054] \\
Feature MMD & 822 & 0.069 & [0.064, 0.073] \\
KS rejection rate & 342 & 0.029 & [0.026, 0.032] \\
Score JSD & 668 & 0.056 & [0.052, 0.060] \\
Prediction-rate shift & 520 & 0.043 & [0.040, 0.047] \\
Prediction JSD & 586 & 0.049 & [0.045, 0.053] \\
Label prior shift & 570 & 0.048 & [0.044, 0.051] \\
Label JSD & 570 & 0.048 & [0.044, 0.051] \\
Confusion profile $L_1$ & 517 & 0.043 & [0.040, 0.047] \\
Confusion profile JSD & 638 & 0.053 & [0.049, 0.057] \\
\midrule
$\mathcal I_X$ (any sensor) & 1500 & 0.125 & [0.119, 0.131] \\
$\mathcal I_{XF}$ (any sensor) & 2438 & 0.203 & [0.196, 0.210] \\
$\mathcal I_{Y_m}$ (any sensor) & 570 & 0.048 & [0.044, 0.051] \\
$\mathcal I_{XFY}$ (any sensor) & 3113 & 0.259 & [0.252, 0.267] \\
\bottomrule
\end{tabular}
\end{table}

\begin{table*}[!t]
\caption{Gate N: per-environment false-alarm rates on 1,200 evaluation draws (three dimensions, two or four models, 200 draws per cell). Sensor range is the minimum and maximum per-sensor rate in the environment; regime columns are batch-level unions. E1 draws 256 rows from 322-row pool halves and is indicative only.}
\label{tab:s-fpr-env}
\centering
\footnotesize
\setlength{\tabcolsep}{5pt}
\begin{tabular}{@{}lrrlrrrr@{}}
\toprule
Environment & $n$ & Pool half & Sensor range & $\mathcal I_X$ & $\mathcal I_{XF}$ & $\mathcal I_{Y_m}$ & $\mathcal I_{XFY}$ \\
\midrule
E1 CICIDS ID 256 & 256 & 322 & 0.023--0.130 & 0.203 & 0.245 & 0.045 & 0.331 \\
E2 UNSW-NB15 ID & 128 & 386 & 0.008--0.070 & 0.063 & 0.135 & 0.070 & 0.210 \\
E3 ToN-IoT ID & 128 & 386 & 0.025--0.172 & 0.215 & 0.280 & 0.070 & 0.358 \\
E4 Tue$\to$Wed & 128 & 1436 & 0.025--0.082 & 0.150 & 0.230 & 0.025 & 0.254 \\
E5 Tue$\to$Fri-PortScan & 128 & 1436 & 0.018--0.051 & 0.066 & 0.160 & 0.025 & 0.198 \\
E6 Wed$\to$Thu-Web & 128 & 1436 & 0.013--0.057 & 0.070 & 0.163 & 0.025 & 0.181 \\
E7 Wed$\to$Fri-AM & 128 & 1436 & 0.025--0.066 & 0.133 & 0.212 & 0.025 & 0.258 \\
E8 UNSW temporal OOD & 128 & 1436 & 0.032--0.092 & 0.072 & 0.192 & 0.050 & 0.237 \\
\bottomrule
\end{tabular}
\end{table*}

\begin{table*}[!t]
\caption{Gate N: detection rate of each regime on the frozen \texttt{paper\_core} observations (all environments, dimensions and models pooled; 600 observations per intervention row) at the null-calibrated thresholds, and firing rate on the in-place near-null controls. Batch-level regimes use calibrated distributional sensors; item-aligned regimes use exact zero thresholds on item-wise deltas.}
\label{tab:s-detection}
\centering
\footnotesize
\setlength{\tabcolsep}{5pt}
\begin{tabular}{@{}lrrrrrr@{}}
\toprule
Intervention & $\mathcal I_X$ & $\mathcal I_{XF}$ & $\mathcal I_{Y_m}$ & $\mathcal I_{XFY}$ & $\mathcal I_{XF}$ item & $\mathcal I_{XFY}$ item \\
\midrule
Feature sign flip 0.02 & 0.118 & 0.192 & 0.000 & 0.200 & 0.793 & 0.793 \\
Feature sign flip 0.05 & 0.532 & 0.640 & 0.000 & 0.643 & 0.867 & 0.867 \\
Feature sign flip 0.10 & 0.757 & 0.877 & 0.000 & 0.878 & 0.907 & 0.907 \\
Mean shift 0.02 & 1.000 & 1.000 & 0.000 & 1.000 & 0.357 & 0.357 \\
Mean shift 0.05 & 1.000 & 1.000 & 0.000 & 1.000 & 0.500 & 0.500 \\
Mean shift 0.10 & 1.000 & 1.000 & 0.000 & 1.000 & 0.618 & 0.618 \\
Scaling drift 0.02 & 1.000 & 1.000 & 0.000 & 1.000 & 0.537 & 0.537 \\
Scaling drift 0.05 & 1.000 & 1.000 & 0.000 & 1.000 & 0.608 & 0.608 \\
Scaling drift 0.10 & 1.000 & 1.000 & 0.000 & 1.000 & 0.668 & 0.668 \\
Feature dropout 0.02 & 0.000 & 0.000 & 0.000 & 0.000 & 0.488 & 0.488 \\
Feature dropout 0.05 & 0.000 & 0.022 & 0.000 & 0.022 & 0.705 & 0.705 \\
Feature dropout 0.10 & 0.000 & 0.062 & 0.000 & 0.062 & 0.763 & 0.763 \\
Prior-preserving label flip 0.02 & 0.000 & 0.000 & 0.000 & 0.000 & 0.000 & 1.000 \\
Prior-preserving label flip 0.05 & 0.000 & 0.000 & 0.000 & 0.002 & 0.000 & 1.000 \\
Prior-preserving label flip 0.10 & 0.000 & 0.000 & 0.000 & 0.033 & 0.000 & 1.000 \\
Random label flip 0.02 & 0.000 & 0.000 & 0.000 & 0.000 & 0.000 & 1.000 \\
Random label flip 0.05 & 0.000 & 0.000 & 0.000 & 0.003 & 0.000 & 1.000 \\
Random label flip 0.10 & 0.000 & 0.000 & 0.000 & 0.033 & 0.000 & 1.000 \\
\midrule
Near-null Gaussian $\sigma=0.001$ & 0.713 & 0.713 & 0.000 & 0.713 & 0.042 & 0.042 \\
Near-null scaling $\alpha=0.001$ & 0.879 & 0.882 & 0.000 & 0.882 & 0.099 & 0.099 \\
\bottomrule
\end{tabular}
\end{table*}

\begin{table}[!t]
\caption{Gate N: label-path interventions in the frozen expansion; number of observations on which each regime fires after calibration (batch-level) or by exact item-aligned comparison.}
\label{tab:s-labelpath}
\centering
\scriptsize
\setlength{\tabcolsep}{2.5pt}
\begin{tabular}{@{}lrrrrrr@{}}
\toprule
Subset & $n$ & $\mathcal I_X$ & $\mathcal I_{XF}$ & $\mathcal I_{Y_m}$ & $\mathcal I_{XFY}$ & item-aligned \\
\midrule
All label interventions & 3,600 & 0 & 0 & 0 & 43 & 3,600 \\
Prior-preserving & 1,800 & 0 & 0 & 0 & 21 & 1,800 \\
Positive conclusion impact & 2,184 & 0 & 0 & 0 & 40 & 2,184 \\
\bottomrule
\end{tabular}
\end{table}

\begin{table}[!t]
\caption{Cluster structure of the frozen evaluation batches (mean over the 30 cells of each environment, standardized projected features): fraction of near-duplicate rows, and the largest fraction of rows within a window of 0.01 or 0.001 standard deviations in any feature (max) or on average over features (mean).}
\label{tab:s-mass}
\centering
\scriptsize
\setlength{\tabcolsep}{3pt}
\begin{tabular}{@{}lrrrr@{}}
\toprule
Environment & Dup.\ rows & Max 0.01 & Mean 0.01 & Max 0.001 \\
\midrule
E1 CICIDS ID 256 & 0.074 & 0.825 & 0.562 & 0.648 \\
E2 UNSW-NB15 ID & 0.114 & 0.735 & 0.458 & 0.514 \\
E3 ToN-IoT ID & 0.257 & 0.870 & 0.466 & 0.571 \\
E4 Tue$\to$Wed & 0.054 & 0.766 & 0.489 & 0.480 \\
E5 Tue$\to$Fri-PortScan & 0.129 & 0.872 & 0.747 & 0.822 \\
E6 Wed$\to$Thu-Web & 0.004 & 0.827 & 0.433 & 0.696 \\
E7 Wed$\to$Fri-AM & 0.026 & 0.769 & 0.670 & 0.663 \\
E8 UNSW temporal OOD & 0.100 & 0.694 & 0.437 & 0.477 \\
\bottomrule
\end{tabular}
\end{table}


\subsection{Gate P: Symmetric Preprocessing Ablation}

The CICIDS ID design (128/128, \texttt{paper\_core}, ZZ and SVC, five splits,
two model seeds, three dimensions) is re-executed under three scaler
configurations: P0, the frozen reference (\texttt{standard} for SVC,
\texttt{minmax2$\pi$} for QSVC); PA, \texttt{standard} for both; PB,
\texttt{minmax2$\pi$} for both. The endpoint is the within-split ZZ-minus-SVC
difference in suite-average balanced-accuracy conclusion impact with a
95\% $t$ interval over the five split clusters; clean balanced accuracy and
the exact label-path invariance counts are reported for every configuration.

\subsection{Gate T: Cross-Validated Tuning of Both Learners}

The same design is re-executed with both learners tuned on training rows only:
SVC over $C\in\{0.1,1,10,100\}\times\gamma\in\{\text{scale},0.01,0.1,1\}$ and
QSVC over $C\in\{0.1,1,10,100\}$ on the precomputed training fidelity kernel,
with stratified five-fold cross-validation (shuffle, random state = model
seed), balanced-accuracy scoring, and ties resolved toward the smallest $C$.
Two environments were fixed before execution: CICIDS ID and the expansion OOD
environment with the highest mean clean SVC balanced accuracy in the frozen
1.0.0 tables, which is UNSW temporal OOD (E8, 0.753; the CICIDS temporal pairs
lie between 0.485 and 0.493). Untuned references with identical seeds are P0
for CICIDS ID and the frozen E8 rows for UNSW OOD.

\textbf{Results of Gates P and T.} Table~\ref{tab:s-paired} reports the
paired differences and Table~\ref{tab:s-clean} the clean balanced accuracies
and selected regularisation constants. The direction of the secondary profile
is stable (five of five positive split clusters in every configuration,
dimension and environment), but its magnitude is conditional: standardizing
the quantum branch (PA) reduces the difference to roughly a third of the
reference, matched min--max scaling (PB) leaves it unchanged, tuning both
learners widens it in CICIDS ID as ZZ gains clean headroom, and leaves UNSW
temporal OOD unchanged while training-only tuning slightly lowers the OOD
accuracy of SVC. Exact label-path invariance holds in every configuration
(360/360 label rows blind to feature/prediction evidence, 180/180
prior-preserving rows blind to label marginals, every positive impact with
joint response, no negative impact).

\begin{table*}[!t]
\caption{Gates P and T: within-split ZZ-minus-SVC difference in suite-average balanced-accuracy conclusion impact, mean with 95\% $t$ interval over five split clusters and number of positive clusters. CICIDS ID reference is the re-executed P0 configuration; UNSW OOD reference is the frozen 1.0.0 environment E8.}
\label{tab:s-paired}
\centering
\footnotesize
\setlength{\tabcolsep}{4pt}
\begin{tabular}{@{}llll@{}}
\toprule
Condition & $d=8$ & $d=10$ & $d=12$ \\
\midrule
P0 reference (standard / minmax2$\pi$) & 0.0245 [0.0069, 0.0422] (5/5) & 0.0348 [0.0176, 0.0520] (5/5) & 0.0407 [0.0305, 0.0509] (5/5) \\
PA standard / standard & 0.0070 [0.0014, 0.0125] (5/5) & 0.0130 [0.0041, 0.0219] (5/5) & 0.0111 [0.0052, 0.0170] (5/5) \\
PB minmax2$\pi$ / minmax2$\pi$ & 0.0238 [0.0054, 0.0423] (5/5) & 0.0323 [0.0156, 0.0491] (5/5) & 0.0396 [0.0261, 0.0531] (5/5) \\
CICIDS ID, both tuned (CV5) & 0.0426 [0.0094, 0.0757] (5/5) & 0.0713 [0.0494, 0.0932] (5/5) & 0.0776 [0.0610, 0.0943] (5/5) \\
UNSW OOD, frozen defaults & 0.0226 [0.0148, 0.0304] (5/5) & 0.0379 [0.0260, 0.0498] (5/5) & 0.0483 [0.0400, 0.0565] (5/5) \\
UNSW OOD, both tuned (CV5) & 0.0229 [0.0103, 0.0356] (5/5) & 0.0348 [0.0205, 0.0492] (5/5) & 0.0459 [0.0371, 0.0546] (5/5) \\
\bottomrule
\end{tabular}
\end{table*}

\begin{table*}[!t]
\caption{Gates P and T: clean balanced accuracy (SVC / ZZ, mean over five split clusters) per configuration and dimension, and the distribution of cross-validated regularisation constants over the 30 cells of each tuned environment (SVC also selects $\gamma$; see the artifact table).}
\label{tab:s-clean}
\centering
\scriptsize
\setlength{\tabcolsep}{3pt}
\begin{tabular}{@{}llll@{}}
\toprule
Condition & $d=8$ & $d=10$ & $d=12$ \\
\midrule
P0 reference (standard / minmax2$\pi$) & 0.753 / 0.734 & 0.747 / 0.728 & 0.778 / 0.747 \\
PA standard / standard & 0.753 / 0.727 & 0.747 / 0.720 & 0.778 / 0.727 \\
PB minmax2$\pi$ / minmax2$\pi$ & 0.751 / 0.734 & 0.750 / 0.728 & 0.756 / 0.747 \\
CICIDS ID, both tuned (CV5) & 0.741 / 0.750 & 0.746 / 0.781 & 0.780 / 0.802 \\
UNSW OOD, frozen defaults & 0.748 / 0.760 & 0.754 / 0.766 & 0.755 / 0.769 \\
UNSW OOD, both tuned (CV5) & 0.739 / 0.759 & 0.723 / 0.757 & 0.736 / 0.770 \\
\midrule
Tuned environment & \multicolumn{2}{l}{QSVC $C$ selections} & SVC $C$ selections \\
\midrule
CICIDS ID & \multicolumn{2}{l}{$C=0.1$: 3, $C=1$: 4, $C=10$: 6, $C=100$: 17} & $C=0.1$: 1, $C=1$: 6, $C=10$: 15, $C=100$: 8 \\
UNSW OOD & \multicolumn{2}{l}{$C=1$: 17, $C=10$: 13} & $C=0.1$: 1, $C=1$: 2, $C=10$: 15, $C=100$: 12 \\
\bottomrule
\end{tabular}
\end{table*}


\section{Reproduction and Evidence Manifest}

The recommended read-only verification is:
{\scriptsize
\begin{verbatim}
python -m src.experiments.`
verify_publication_artifact `
--root publication/artifact
\end{verbatim}
}
It verifies the artifact-wide SHA-256 manifest, five embedded evidence
manifests, 44 outputs, row counts, primary-claim counts, and the calibrated
label-path consistency checks of the reinforcement gate. In artifact 1.1.0 it
reports 3,600 expansion label rows, 2,184 positive expansion impacts, 1,440
Gate-1 label rows, 1,276 positive Gate-1 impacts, and 60 calibrated
(environment, dimension, model) cells. The 26-test suite covers invariances,
the exact kernel, the contracts, the null-control sampler, the frozen
configuration fingerprint, the prespecified threshold rank, and the
cross-validation tie rule.

The complete test command is:
{\scriptsize
\begin{verbatim}
python -m pytest tests -q -p no:cacheprovider
\end{verbatim}
}
The frozen environment uses Python 3.10.16, NumPy 2.2.6, pandas 2.3.3,
scikit-learn 1.7.2, Qiskit 2.3.0, and Qiskit Machine Learning 0.9.0. Qiskit is
pinned below 3 because migrating deprecated feature-map constructors would
change circuit serialization and therefore requires a new evidence version.

The artifact intentionally omits raw benchmark datasets. It includes staging
instructions, source locations, source code, locked dependencies, compact
derived evidence, figures, figure-source tables, tests, manifests, and expected
outputs. Version 1.1.1 is archived at Zenodo, DOI 10.5281/zenodo.22552643,
with the repository tag \texttt{paper15-q1-v1.1.1} (concept DOI
10.5281/zenodo.22550852; version 1.1.0, DOI 10.5281/zenodo.22550853, differs
only in the wording of the funding acknowledgement); source code is licensed
under Apache-2.0, derived evidence, figures, and documentation under CC BY 4.0,
and the manuscript files are author preprints excluded from both licenses.

\section{Claims-Supported / Claims-Excluded Matrix}

\textbf{Supported by the artifact.}
Auditability is conditional on evidence exposed at a declared boundary;
model-output and evaluation-conclusion impacts can have different causal loci;
feature/prediction sensors are exactly blind to label-only changes, whereas
label marginals are blind to prior-preserving relabeling; these blind regions
persist under calibration against a nominal $\alpha=0.05$ per-sensor target and
are closed only by item-aligned evidence; the non-invariant sensors, calibrated
against that target, show pooled empirical false-alarm rates of 0.029--0.069 on
disjoint clean pools within the frozen design, with uncorrected regime unions
of 0.125--0.259; the secondary ZZ-minus-SVC
profile is direction-stable but scaler- and headroom-dependent; complementary
quantum-workflow sensors close different simulator-gate blind regions; and the
local prototype executes policy-aware, fail-closed
\texttt{allow/hold/block} composition.

\textbf{Explicitly excluded.}
The evidence does not estimate attack prevalence or population-level deployment
incidence, establish universal superiority or vulnerability of quantum
kernels, or causally isolate feature-map geometry from preprocessing and clean
headroom. It provides no evidence about QPU security, calibrated device noise,
scheduling, provider integrity, multi-tenancy, side channels, production
readiness, provider authentication, non-repudiation, persistence, incident
handling, or operational Fleet Management.

\end{document}
